\documentclass[conference]{IEEEtran}
\IEEEoverridecommandlockouts
\usepackage{cite}
\usepackage{amsmath,amssymb,amsfonts}
\usepackage{algorithmic}
\usepackage{graphicx}
\usepackage{textcomp}
\usepackage{xcolor}
\usepackage{booktabs}
\usepackage{verbatim}
\def\BibTeX{{\rm B\kern-.05em{\sc i\kern-.025em b}\kern-.08em
    T\kern-.1667em\lower.7ex\hbox{E}\kern-.125emX}}
\begin{document}

\title{Lista de Exercícios: Máquinas de Turing, Decidibilidade e Complexidade\\
{\footnotesize Fundamentos da Computação}
}

\author{\IEEEauthorblockN{Amilton Rocha Holanda}
\IEEEauthorblockA{\textit{Universidade de Fortaleza - UNIFOR}\\
Ciências da Computação\\
Matrícula: 2026169}
\and
\IEEEauthorblockN{Isaac Newton Montinegro}
\IEEEauthorblockA{\textit{Universidade de Fortaleza - UNIFOR}\\
Ciências da Computação\\
Matrícula: 2320324}
}

\maketitle

\begin{abstract}
Este trabalho apresenta a resolução de uma lista de exercícios sobre Máquinas de Turing, decidibilidade e complexidade computacional. São abordados conceitos fundamentais como a Tese de Church-Turing, o problema da parada, as classes P, NP, NP-completo e NP-difícil, além da implementação de Máquinas de Turing para problemas específicos. O trabalho também discute as implicações práticas da teoria da complexidade na criptografia, com ênfase no RSA.
\end{abstract}

\begin{IEEEkeywords}
Máquina de Turing, Decidibilidade, Complexidade Computacional, P vs NP, RSA, Church-Turing
\end{IEEEkeywords}

\section{Área 1: Fundamentos da Computação e da Máquina de Turing}

\subsection{Questão 1}
\textbf{Como a Tese de Church-Turing contribui para os limites da computação moderna?}

A Tese de Church-Turing estabelece que qualquer função computável pode ser computada por uma Máquina de Turing. Isso define o limite do que é computável: se um problema não pode ser resolvido por uma MT, ele não pode ser resolvido por \textit{nenhum} algoritmo computacional. Na computação moderna, isso significa que problemas como o Problema da Parada são intrinsecamente insolúveis, independentemente do hardware ou software disponível. A tese também fundamenta a teoria da computabilidade e orienta o desenvolvimento de linguagens de programação, que são equivalentes a MTs (Turing-completas).

\subsection{Questão 2}
\textbf{Por que o problema da parada é indecidível e o que isso implica para a computação?}

O problema da parada pergunta se existe um algoritmo que, dada uma descrição de um programa e sua entrada, determine se o programa irá parar ou executar infinitamente. Turing provou sua indecidibilidade por contradição: suponha que exista um programa $H(P, I)$ que decide se $P$ para com entrada $I$. Construímos um programa $D(P)$ que chama $H(P, P)$ e entra em loop se $H$ disser que para, e para se $H$ disser que não para. Então $D(D)$ leva a uma contradição. Isso implica que existem limites fundamentais para a computação: não é possível criar um depurador universal que detecte loops infinitos, e muitos problemas interessantes em verificação de software são indecidíveis.

\subsection{Questão 3}
\textbf{A fita infinita é uma abstração realista? Por que ou por que não? É útil tê-la de modo infinito?}

A fita infinita não é realista do ponto de vista físico, pois recursos de memória são sempre finitos. No entanto, é uma abstração matemática necessária para modelar a computação sem limites artificiais. Ela é útil porque permite raciocinar sobre o que é computável em princípio, sem restrições de memória que seriam arbitrárias. Na prática, qualquer computação real usa uma quantidade finita de memória, mas considerar a fita infinita permite separar problemas de computabilidade (o que é possível computar) de problemas de complexidade (com que recursos).

\subsection{Questão 4}
\textbf{A ideia de que a MT não resolve todos os problemas significa que ela é limitada?}

Não. O fato de a MT não resolver todos os problemas não é uma limitação da MT, mas sim uma propriedade fundamental da computação. A Tese de Church-Turing afirma que a MT captura exatamente a noção de computabilidade algorítmica. Portanto, se um problema não pode ser resolvido por uma MT, ele é intrinsecamente não algorítmico — nenhum computador, presente ou futuro, poderá resolvê-lo. A MT não é limitada; ela define o próprio limite da computação.

\subsection{Questão 5}
\textbf{Como o conceito de computabilidade se relaciona com decidibilidade?}

Computabilidade refere-se à capacidade de uma função ser computada por uma MT, enquanto decidibilidade refere-se à capacidade de um problema de decisão (resposta sim/não) ser resolvido por uma MT que sempre para. Um problema de decisão é decidível se sua função característica é computável. Portanto, decidibilidade é um caso particular de computabilidade: problemas decidíveis são aqueles cuja resposta pode ser computada para toda entrada. Problemas indecidíveis são aqueles cuja função característica não é computável.

\subsection{Questão 6}
\textbf{Por que a tese de Church-Turing é um alicerce fundamental para a teoria da computação?}

A Tese de Church-Turing é fundamental porque unifica diferentes formalismos (cálculo lambda, funções recursivas, máquinas de Turing) em uma única noção de computabilidade. Ela permite que a ciência da computação estude os limites fundamentais da computação de forma matemática e rigorosa. Além disso, serve como base para a teoria da complexidade, para a computabilidade e para o projeto de linguagens de programação (que são Turing-completas). Sem ela, não haveria um padrão para determinar se um problema é ou não computável.

\subsection{Questão 7}
\textbf{Todas as variantes da máquina de Turing têm o mesmo poder computacional? Explique.}

Sim, todas as variantes da MT (multifita, não determinística, multidimensional, com múltiplos cabeçotes, entre outras) têm o mesmo poder computacional que a MT padrão de fita única. Isso significa que qualquer problema resolvido por uma variante pode ser simulado pela MT padrão. A diferença está na eficiência: uma MT multifita pode resolver um problema em menos passos que uma MT de fita única, mas ambas podem resolver exatamente a mesma classe de problemas. Isso demonstra a robustez do modelo de Turing.

\subsection{Questão 8}
\textbf{Qual a implicação prática de um problema ser indecidível?}

A implicação prática é que não existe algoritmo geral que resolva o problema para todas as entradas possíveis. Por exemplo, não é possível criar um programa antivírus que detecte perfeitamente todos os vírus (isso é equivalente ao problema da parada). Na engenharia de software, não é possível verificar automaticamente se um programa qualquer termina ou se comporta corretamente em todos os casos. Isso leva ao uso de aproximações, heurísticas e análise estática limitada, que funcionam para muitos casos mas não para todos.

\subsection{Questão 9}
\textbf{Como a redutibilidade ajuda a demonstrar a indecidibilidade de problemas?}

A redutibilidade transforma um problema $A$ em um problema $B$ de modo que uma solução para $B$ implicaria uma solução para $A$. Se $A$ é indecidível e $A$ é redutível a $B$, então $B$ também é indecidível (caso contrário, poderíamos resolver $A$ resolvendo $B$). Por exemplo, reduzimos o problema da parada a vários problemas indecidíveis. Isso cria uma cadeia de indecidibilidade: uma vez que se conhece um problema indecidível (como o da parada), pode-se demonstrar que outros são indecidíveis por redução.

\subsection{Questão 10}
\textbf{Uma linguagem é considerada Turing-decidível quando:}

As alternativas são:
A) Uma Máquina de Turing entra em loop para todas as entradas. \\
B) Toda entrada válida é aceita pela Máquina de Turing. \\
C) A Máquina de Turing sempre para em estados de aceitação ou rejeição. \\
D) Ela pertence à classe P. \\
E) Não pode ser reconhecida por uma Máquina de Turing.

\textbf{Resposta: Alternativa C.} Uma linguagem é Turing-decidível (ou recursiva) quando existe uma MT que sempre para — seja em estado de aceitação (se a cadeia pertence à linguagem) ou de rejeição (se não pertence). As alternativas A, B e E descrevem MTs que não decidem (loop, aceitação parcial, não reconhecimento). A alternativa D confunde decidibilidade com complexidade (classe P);

\subsection{Questão 12}
\textbf{(Concurso IF-SC - 2014) Assinale a alternativa CORRETA:}

A) A máquina de Turing consiste de uma fita finita; um cabeçote que lê, escreve e move para direita ou esquerda; um registrador de estados e uma tabela de ações. \\
B) O problema da parada da máquina de Turing deve-se ao limite finito de sua fita e as poucas operações que um cabeçote pode executar (avançar ou retroceder). \\
C) A máquina de Turing pode ser considerada um autômato infinito de grau dois. \\
D) Se um problema não puder ser resolvido por uma máquina de Turing, então esse problema não poderá ser resolvido por qualquer outro sistema algorítmico. \\
E) Criar uma máquina de Turing com fita infinita ainda não é possível devido às restrições tecnológicas atuais.

\textbf{Resposta: Alternativa D.} Esta é a essência da Tese de Church-Turing: se um problema não pode ser resolvido por uma MT, não pode ser resolvido por nenhum algoritmo. A alternativa A está errada porque a fita é infinita (não finita). B está errada porque o problema da parada não se deve a limites físicos. C está errada pois MT não é um autômato infinito de grau dois. E está errada pois a fita é uma abstração matemática, não uma restrição tecnológica.

\section{Área 2: Desenvolvimento de Máquina de Turing}

\subsection{Questão 13/14 — Multiplicação em Notação Unária}

Implementar uma Máquina de Turing que valida a multiplicação de dois números naturais em notação unária. Na fita temos o formato $1^a X 1^b = 1^{a \times b}$. Por exemplo, para $3 \times 2 = 6$, a fita contém \texttt{111X11=111111}. A MT deve verificar se a quantidade de 1s após o sinal de igual é o produto das quantidades antes do X e do =.

\subsubsection{Descrição do Algoritmo}

\begin{enumerate}
  \item A MT inicia no símbolo mais à esquerda e percorre os 1s do primeiro fator, substituindo cada 1 por \texttt{\_\_\_} (underscore) e copiando o segundo fator para o final da fita, para cada 1 do primeiro fator.
  \item Ao encontrar o X, a MT marca a posição e começa a percorrer o segundo fator.
  \item Para cada 1 do primeiro fator, a MT copia todo o bloco do segundo fator para o final da fita (após o =).
  \item Ao final, compara a quantidade de 1s após o = com o produto esperado. Na versão de validação, a MT simplesmente verifica se a quantidade de 1s é igual ao produto.
\end{enumerate}

\subsubsection{Implementação em Código}

A implementação em Python que simula a MT para validação da multiplicação unária está no arquivo \texttt{mt\_multiplicacao.py}. O código foi estruturado com:
\begin{itemize}
  \item Estados nomeados para cada etapa do processo
  \item Função de transição que mapeia (estado, símbolo) $\rightarrow$ (novo estado, símbolo escrito, direção)
  \item Simulação passo a passo com registro da fita a cada transição
\end{itemize}

\subsection{Questão 16 — Cópia de Sequência Binária}

Implementar uma Máquina de Turing que copie uma sequência binária separada por um marcador especial \#. Entrada: $101\#$, saída esperada: $101\#101$.

\subsubsection{Descrição do Algoritmo}

\begin{enumerate}
  \item A MT lê cada símbolo da sequência antes do \# e escreve o símbolo em uma posição temporária ou marca o símbolo.
  \item Ao chegar no \#, a MT retorna ao início da fita.
  \item Para cada símbolo lido (da esquerda para a direita até o \#), a MT o copia para após o \#.
  \item O processo termina quando todos os símbolos foram copiados.
\end{enumerate}

Uma estratégia eficiente:
\begin{itemize}
  \item Substituir cada 0 por \texttt{X} e cada 1 por \texttt{Y} para marcar que já foi copiado.
  \item Para cada símbolo marcado, ir até o final da fita (após \#) e escrever o 0 ou 1 correspondente.
  \item Retornar ao início e repetir até que todos os símbolos antes de \# sejam marcados.
  \item Finalmente, restaurar os símbolos originais na primeira parte.
\end{itemize}

\subsubsection{Implementação em Código}

A implementação em Python está no arquivo \texttt{mt\_copia\_binaria.py}.

\section{Área 3: Complexidade Computacional}

\subsection{Questão 17}
\textbf{A classe P é sempre eficiente na prática?}

Não necessariamente. Um algoritmo $O(n^{100})$ está em P, mas é completamente impraticável. Na prática, algoritmos polinomiais com expoentes pequenos ($n$, $n^2$, $n^3$) são eficientes, mas expoentes maiores podem ser tão ruins quanto exponenciais para entradas modestas. A classe P é uma definição teórica que captura a noção de eficiência assintótica, mas a eficiência prática depende do tamanho do expoente e das constantes envolvidas.

\subsection{Questão 18}
\textbf{Por que problemas da classe NP são considerados desafiadores?}

Porque não se conhece algoritmos polinomiais para resolvê-los, apenas para verificar soluções. Suspeita-se que P $\neq$ NP, o que significaria que esses problemas são intrinsecamente difíceis de resolver. Além disso, muitos problemas NP têm importância prática (roteamento, escalonamento, otimização) e mesmo instâncias moderadas podem exigir tempo exponencial.

\subsection{Questão 19}
\textbf{Qual a diferença entre resolver e verificar uma solução em termos de P e NP?}

Resolver significa encontrar a solução a partir da entrada. Verificar significa, dada uma solução candidata (certificado), confirmar se ela é correta. P é a classe dos problemas que podem ser resolvidos em tempo polinomial. NP é a classe dos problemas cujas soluções podem ser verificadas em tempo polinomial. Por exemplo, fatorar um número é difícil (resolver), mas verificar se uma fatoração está correta é trivial (multiplicar).

\subsection{Questão 20}
\textbf{Se um problema NP-completo for resolvido em tempo polinomial, o que acontece com os demais problemas em NP?}

Se um problema NP-completo for resolvido em tempo polinomial, então P = NP. Isso ocorre porque problemas NP-completos são os "mais difíceis" em NP — qualquer problema em NP pode ser reduzido em tempo polinomial a um problema NP-completo. Se o NP-completo tem solução polinomial, então todos os problemas em NP também têm.

\subsection{Questão 21}
\textbf{Por que a questão P = NP é considerada um dos maiores problemas abertos da ciência?}

Porque está no centro da computação teórica e tem implicações profundas em criptografia, otimização, inteligência artificial e biologia computacional. Se P = NP, problemas considerados intratáveis (como otimização de rotas, dobramento de proteínas, etc.) teriam soluções eficientes. Se P $\neq$ NP, confirmar-se-ia que certos problemas são fundamentalmente difíceis. O Clay Mathematics Institute oferece US\$ 1 milhão pela solução.

\subsection{Questão 22}
\textbf{O que significaria, para a criptografia, provar que P = NP?}

Seria catastrófico para a criptografia moderna. A segurança de sistemas como RSA depende da dificuldade de fatorar números grandes (problema em NP). Se P = NP, existiria um algoritmo polinomial para fatoração, quebrando o RSA e outros criptossistemas de chave pública. Toda a comunicação segura na internet (HTTPS, SSH, etc.) seria comprometida. Novos paradigmas seriam necessários.

\subsection{Questão 23}
\textbf{Por que nem todos os problemas NP são NP-completos?}

NP-completo é uma subclasse de NP que contém os problemas "mais difíceis" de NP. Para um problema ser NP-completo, todo problema em NP deve ser redutível a ele. Existem problemas em NP que não são NP-completos, como problemas em P (que estão em NP mas são fáceis) e problemas NP-intermediários (como fatoração e isomorfismo de grafos), para os quais não se sabe se estão em P nem se são NP-completos.

\subsection{Questão 24}
\textbf{Como se diferencia um problema NP-completo de um NP-difícil?}

NP-completo é a interseção de NP e NP-difícil: o problema deve estar em NP (verificável em tempo polinomial) e ser NP-difícil (todo problema em NP é redutível a ele). NP-difícil é uma classe mais ampla: problemas que são pelo menos tão difíceis quanto qualquer problema em NP, mas que podem não estar em NP (podem ser indecidíveis ou exigir mais que tempo polinomial para verificação).

\subsection{Questão 25}
\textbf{Por que a verificabilidade em tempo polinomial é uma propriedade crucial?}

Porque reflete a intuição de que reconhecer uma solução correta é mais fácil que encontrá-la. Muitos problemas do mundo real têm essa propriedade: dado um candidato a solução, é fácil verificar se é válido, mas difícil de encontrar. A verificabilidade polinomial define precisamente a classe NP, que captura uma grande variedade de problemas práticos. Essa propriedade também é base para a criptografia de chave pública.

\subsection{Questão 26}
\textbf{Análise de conceitos de P e NP:}

Afirmativas:
I. Todo problema em P está também em NP. \\
II. Problemas NP têm soluções verificáveis em tempo polinomial. \\
III. Se um problema NP-completo for resolvido em tempo polinomial, então P = NP.

\textbf{Resposta: Alternativa C (I, II e III).} 
I é verdadeira: se um problema pode ser resolvido em tempo polinomial, sua solução pode ser verificada em tempo polinomial. 
II é verdadeira: esta é a definição de NP (verificável em tempo polinomial). 
III é verdadeira: pela definição de NP-completude, todo problema em NP se reduz a um NP-completo.

\subsection{Questão 27}
\textbf{Sobre limites computacionais, a indecidibilidade é característica de problemas que:}

A) Têm solução em tempo exponencial. \\
B) Não podem ser resolvidos por nenhuma máquina de Turing. \\
C) Pertencem à classe P. \\
D) Podem ser sempre resolvidos por algoritmos não determinísticos.

\textbf{Resposta: Alternativa B.} Problemas indecidíveis não podem ser resolvidos por nenhuma MT (e portanto por nenhum algoritmo).

\subsection{Questão 28}
\textbf{A redutibilidade entre problemas é essencial para demonstrar indecidibilidade porque:}

A) Permite transformar problemas indecidíveis em decidíveis. \\
B) Facilita a implementação de algoritmos mais rápidos. \\
C) Estabelece equivalência de complexidade entre problemas distintos. \\
D) Garante que todo problema NP é também NP-completo.

\textbf{Resposta: Alternativa C.} A redutibilidade estabelece que se $A \leq B$, então $B$ é pelo menos tão difícil quanto $A$. Para indecidibilidade, reduzir um problema indecidível a $B$ mostra que $B$ também é indecidível.

\subsection{Questão 29}
\textbf{Um problema resolvido por uma MT mas que requer tempo exponencial é:}

A) Indecidível. \\
B) NP-completo. \\
C) Intratável. \\
D) Não verificável.

\textbf{Resposta: Alternativa C.} Um problema é intratável se requer recursos exponenciais no pior caso, mesmo sendo decidível.

\subsection{Questão 30}
\textbf{Uma linguagem é Turing-decidível quando:}

F) Uma MT entra em loop para todas as entradas. \\
G) Toda entrada válida é aceita pela MT. \\
H) A MT sempre para em estados de aceitação ou rejeição. \\
I) Ela pertence à classe P. \\
J) Não pode ser reconhecida por uma MT.

\textbf{Resposta: Alternativa H.} Uma linguagem é decidível se existe uma MT que sempre para, respondendo sim (aceita) ou não (rejeita) para qualquer entrada.

\subsection{Questão 31}
\textbf{(ENADE, 2011) Análise de assertivas sobre P e NP:}

I. Existem problemas na classe P que não estão na classe NP. \\
II. Se $A$ é eficientemente transformado em $B$ e $B \in P$, então $A \in P$. \\
III. Se P = NP, então um problema NP-completo pode ser solucionado eficientemente. \\
IV. Se P $\neq$ NP, então existem problemas em P que são NP-completos.

\textbf{Resposta: Alternativa D (II e III).}
I é falsa: P $\subseteq$ NP. \\
II é verdadeira: redução polinomial preserva P. \\
III é verdadeira: se P = NP, NP-completo está em P. \\
IV é falsa: se P $\neq$ NP, nenhum problema em P é NP-completo.

\subsection{Questão 32}
\textbf{(Concurso Marinha, 2011) Assinale a opção correta:}

A) P consiste nos problemas que, dado um certificado, é possível verificar a solução em tempo polinomial. \\
B) NP consiste nos problemas que não pertencem a P e não são verificáveis em tempo polinomial. \\
C) A busca binária é NP-completo. \\
D) Um problema em P não está em NP, exceto NP-completos. \\
E) Se um NP-completo tem solução polinomial, todos os problemas em NP têm solução polinomial.

\textbf{Resposta: Alternativa E.} Esta é a definição de NP-completude: se um NP-completo é polinomial, todos os problemas em NP também são. A alternativa A define NP, não P. B está errada (NP inclui P). C está errada (busca binária é P). D está errada (P $\subseteq$ NP).

\subsection{Questão 33}
\textbf{(ENADE, 2005) Julgue os itens:}

I. P $\subseteq$ NP e NP $\subseteq$ P. \\
II. Se P $\neq$ NP, então P $\cap$ NP-Completo $= \emptyset$. \\
III. Se $Q$ é NP-difícil e $Q \in$ NP, então $Q$ é NP-completo. \\
IV. SAT é NP-difícil e NP-Completo. \\
V. Caminho mais curto em grafo não é P.

\textbf{Resposta: Alternativa B (II, III e IV).} \\
I é falsa: NP $\subseteq$ P não é conhecido. \\
II é verdadeira: se P $\neq$ NP, problemas em P não podem ser NP-completos. \\
III é verdadeira: esta é a definição de NP-completo (NP $\cap$ NP-difícil). \\
IV é verdadeira: SAT é NP-completo (Cook-Levin). \\
V é falsa: caminho mais curto é P (Dijkstra, Floyd-Warshall).

\section{Área 4: Criptografia e Aplicações Práticas}

\subsection{Questão 34}
\textbf{Por que o RSA é considerado seguro apesar de sua simplicidade algorítmica?}

O RSA é seguro devido à dificuldade computacional de fatorar números inteiros grandes. O algoritmo em si é simples: requer apenas exponenciação modular e conhecimento do produto de dois primos grandes. No entanto, a segurança vem do fato de que, sem conhecer os fatores primos, não há algoritmo eficiente conhecido para obter a chave privada a partir da chave pública. Sua simplicidade algorítmica não é uma fraqueza; ao contrário, torna o RSA fácil de implementar e analisar.

\subsection{Questão 35}
\textbf{Por que a geração de chaves RSA é rápida, mas sua quebra é extremamente lenta?}

A geração de chaves requer encontrar dois primos grandes (que pode ser feito eficientemente com testes de primalidade probabilísticos como Miller-Rabin) e calcular o produto $n = p \times q$. A quebra requer fatorar $n$, e não se conhece algoritmo polinomial para fatoração de inteiros grandes. O melhor algoritmo clássico para fatoração (General Number Field Sieve) tem complexidade subexponencial, mas não polinomial. Essa assimetria entre geração (fácil) e quebra (difícil) é a base da criptografia de chave pública.

\subsection{Questão 38}
\textbf{O que significa dizer que a segurança do RSA depende de um problema da classe NP?}

O problema da fatoração de inteiros está em NP: dado um candidato a fatoração, é fácil verificar se está correto (basta multiplicar os fatores). No entanto, não se sabe se a fatoração está em P ou se é NP-completa. Acredita-se que a fatoração seja NP-intermediária (nem P, nem NP-completa). A segurança do RSA depende da suposição de que a fatoração é difícil para números grandes. Se P = NP, a fatoração teria solução polinomial e o RSA seria inseguro.

\subsection{Questão 40}
\textbf{Como o estudo do RSA reforça a importância da teoria da complexidade?}

O RSA é um exemplo concreto de como a separação entre resolver e verificar (a base da classe NP) tem aplicações práticas transformadoras. Sem a teoria da complexidade, não haveria fundamento para acreditar na segurança do RSA. O estudo do RSA também motiva questões abertas em complexidade (P vs NP, dureza da fatoração) e mostra como conceitos teóricos abstratos (redutibilidade, classes de complexidade) têm impacto direto na prática. A segurança da internet moderna repousa sobre conjecturas teóricas da ciência da computação.

\begin{thebibliography}{00}
\bibitem{b1} J. E. Hopcroft, R. Motwani, and J. D. Ullman, \textit{Introduction to Automata Theory, Languages, and Computation}, 3rd ed. Boston, MA: Addison-Wesley, 2007.
\bibitem{b2} M. Sipser, \textit{Introduction to the Theory of Computation}, 3rd ed. Boston, MA: Cengage Learning, 2013.
\bibitem{b3} C. H. Papadimitriou, \textit{Computational Complexity}. Reading, MA: Addison-Wesley, 1994.
\bibitem{b4} S. Russell and P. Norvig, \textit{Artificial Intelligence: A Modern Approach}, 3rd ed. Upper Saddle River, NJ: Prentice Hall, 2010.
\bibitem{b5} R. L. Rivest, A. Shamir, and L. Adleman, ``A method for obtaining digital signatures and public-key cryptosystems,'' \textit{Communications of the ACM}, vol. 21, no. 2, pp. 120--126, 1978.
\end{thebibliography}

\end{document}
